% !TITLE 白马篇
% !AUTHOR 曹植
% !DYNASTY 魏晋
% !GENRE 五言古诗
% !ORDER 6

\begin{poem}
白马饰金羁，连翩西北驰。
借问谁家子，幽并游侠儿。
少小去乡邑，扬声沙漠垂。
宿昔秉良弓，楛矢何参差。
控弦破左的，右发摧月支。
仰手接飞猱，俯身散马蹄。
狡捷过猴猿，勇剽若豹螭。
边城多警急，虏骑数迁移。
羽檄从北来，厉马登高堤。
长驱蹈匈奴，左顾凌鲜卑。
弃身锋刃端，性命安可怀？
父母且不顾，何言子与妻！
名编壮士籍，不得中顾私。
捐躯赴国难，视死忽如归。
\end{poem}

\begin{appreciation}
《白马篇》是曹植的代表作之一，写一个边塞游侠少年驰骋沙场、为国捐躯的故事，是建安文学慷慨任气风格的典型体现。

曹植是曹操的第三子，才高八斗，却一生不得志。他年轻时随父出征，亲历过边塞战事，对游侠生活和战争场面有直接的经验。这首诗虽然写的是别人，但字里行间涌动的是他自己的热血和抱负。

白马饰金羁，连翩西北驰——白马配着金色的马笼头，连翩向西北飞驰。开篇就是一个电影般的画面：一匹白马，一个少年，向边疆疾驰而去。借问谁家子，幽并游侠儿——请问这是谁家的少年？是幽州、并州的游侠子弟。幽并之地（今河北、山西一带）自古多豪侠，民风剽悍，尚武好义。

控弦破左的，右发摧月支。仰手接飞猱，俯身散马蹄——拉弓射中左边的靶子，向右发射摧毁了月支。抬手接住飞跳的猿猴，俯身射散马蹄上的目标。这四句写少年的武艺高强，用了四个动作，每一个都精准有力。狡捷过猴猿，勇剽若豹螭——灵活超过猿猴，勇猛好像豹子和螭龙。这些比喻把一个少年英雄的形象刻画得栩栩如生。

羽檄从北来，厉马登高堤——插着羽毛的紧急军书从北方传来，他策马登上高堤。羽檄是古代最紧急的军书，插羽毛表示十万火急。敌人的入侵打破了和平，也给了少年一个施展抱负的机会。

弃身锋刃端，性命安可怀？父母且不顾，何言子与妻！——投身于刀刃之间，怎么还能顾惜性命？父母都顾不上了，何况妻子儿女！这几句是全诗情感的高潮。他不是不知道生命的可贵，不是不爱家人，但在国难当头的时候，个人的一切都必须让位。这种舍生取义的精神，是中国文化中最高贵的品质之一。

捐躯赴国难，视死忽如归——献身于国难，把死亡看作回家一样平常。这是全诗的结句，也是建安文学最响亮的声音。视死忽如归四个字，把死亡从恐惧变成了从容，从悲剧变成了崇高。这种气魄，后世称为建安风骨。
\end{appreciation}
